% GENERATED FILE. Edit canonical lessons, then run npm run generate:books.
% canonical-source-hash: fnv1a64:149b0d5ecb1f7736
% canonical-lessons: RU-C423-zaderzhka, RU-C423-zaderzhivat, RU-C423-otmena, RU-C423-otmenyat, RU-C423-kompensatsiya

\chapter{Delays and cancellations}
\label{ch:ru-delays-and-cancellations}
\hlchaptermodality{423}

\begin{chapteropening}
\textbf{By the end of this chapter:} \emph{I can deal with delays: a delay, being held up, a cancellation, cancelling and compensation.}

Complete the last lesson of chapter 423: I can deal with delays: a delay, being held up, a cancellation, cancelling and compensation.
\end{chapteropening}

\section[\texorpdfstring{\ru{задержка} (zadérzhka)}{задержка (zadérzhka)}]{\ru{задержка} (zadérzhka) — a delay}

\label{lesson:RU-C423-zaderzhka}



\begin{quote}
Before the new one: say the Russian for a terminal, then the Russian for a flight attendant.
\end{quote}

\subsection*{You'll want to know: \ru{задержка}}
\textbf{\ru{задержка}} (\emph{zadérzhka}) — \textquotedblleft{}a delay\textquotedblright{}.

\begin{quote}
\ru{Из}-\ru{за} \ru{тумана} \ru{была} \ru{большая} \ru{задержка}.
\end{quote}

\begin{quote}
\emph{Iz-za tumána bylá bolsháya zadérzhka.}
\end{quote}

\begin{quote}
Because of the fog there was a long delay.
\end{quote}

Feminine. \emph{\ru{Рейс} \ru{с} \ru{задержкой}} — a delayed flight.

\begin{cousinweb}
\emph{\ru{задержка}} is built on \emph{\ru{задержать}} 'to hold back', from \emph{\ru{за}-} and \emph{\ru{держать}} 'to hold'. \emph{\ru{Держать}} is inherited from Proto-Slavic \emph{dĭržati} 'to hold' (EtymDB).
\end{cousinweb}

\begin{grammarlens}[title={What you've built}]
One more word for telling, explaining and describing at B1.
\end{grammarlens}

\subsection*{Your turn}
\begin{itemize}
\raggedright
  \item \emph{Say it:} \emph{zadérzhka}
  \item \emph{Say it:} \emph{zadérzhka}, once more
  \item \emph{Say it:} a sentence of your own with \emph{zadérzhka}
  \item \emph{Recall:} say the Russian for a terminal, then the Russian for a flight attendant, then say \emph{zadérzhka} again
\end{itemize}

\begin{tcolorbox}[breakable,colback=teal!4,colframe=teal!35!black,title={Before you move on}]
What does \ru{задержка} mean? (\textquotedblleft{}A delay\textquotedblright{}.) Say it once more.
\end{tcolorbox}
\section[\texorpdfstring{\ru{задерживать} (zadérzhivat)}{задерживать (zadérzhivat)}]{\ru{задерживать} (zadérzhivat) — to delay; to detain}

\label{lesson:RU-C423-zaderzhivat}



\begin{quote}
Before the new one: say the Russian for a flight attendant, then the Russian for a delay.
\end{quote}

\subsection*{You'll want to know: \ru{задерживать}}
\textbf{\ru{задерживать}} (\emph{zadérzhivat}) — \textquotedblleft{}to delay; to detain\textquotedblright{}.

\begin{quote}
\ru{Снег} \ru{часто} \ru{задерживает} \ru{поезда}.
\end{quote}

\begin{quote}
\emph{Sneg chásto zadérzhivayet poyezdá.}
\end{quote}

\begin{quote}
Snow often holds up trains.
\end{quote}

Imperfective; perfective \emph{\ru{задержать}} (\emph{\ru{задержу}, \ru{задержишь}}). \emph{\ru{Меня} \ru{задержали} \ru{на} \ru{работе}} — I was kept late at work.

\begin{cousinweb}
\emph{\ru{задерживать}} is built from \emph{\ru{за}-} and \emph{-\ru{держивать}}, from \emph{\ru{держать}} 'to hold', which is inherited from Proto-Slavic \emph{dĭržati} (EtymDB).
\end{cousinweb}

\begin{grammarlens}[title={What you've built}]
One more word for telling, explaining and describing at B1.
\end{grammarlens}

\subsection*{Your turn}
\begin{itemize}
\raggedright
  \item \emph{Say it:} \emph{zadérzhivat}
  \item \emph{Say it:} \emph{zadérzhivat}, once more
  \item \emph{Say it:} a sentence of your own with \emph{zadérzhivat}
  \item \emph{Recall:} say the Russian for a flight attendant, then the Russian for a delay, then say \emph{zadérzhivat} again
\end{itemize}

\begin{tcolorbox}[breakable,colback=teal!4,colframe=teal!35!black,title={Before you move on}]
What does \ru{задерживать} mean? (\textquotedblleft{}To delay; to detain\textquotedblright{}.) Say it once more.
\end{tcolorbox}
\section[\texorpdfstring{\ru{отмена} (otména)}{отмена (otména)}]{\ru{отмена} (otména) — a cancellation}

\label{lesson:RU-C423-otmena}



\begin{quote}
Before the new one: say the Russian for a delay, then the Russian for to delay.
\end{quote}

\subsection*{You'll want to know: \ru{отмена}}
\textbf{\ru{отмена}} (\emph{otména}) — \textquotedblleft{}a cancellation\textquotedblright{}.

\begin{quote}
\ru{После} \ru{отмены} \ru{рейса} \ru{мы} \ru{вернулись} \ru{домой}.
\end{quote}

\begin{quote}
\emph{Pósle otmény réysa my vernúlis domóy.}
\end{quote}

\begin{quote}
After the flight was cancelled we went back home.
\end{quote}

Feminine.

\begin{cousinweb}
\emph{\ru{отмена}} is built from \emph{\ru{от}-} 'away' and \emph{\ru{мена}} 'an exchange', which is inherited from Proto-Slavic \emph{měna} 'change' (EtymDB).
\end{cousinweb}

\begin{grammarlens}[title={What you've built}]
One more word for telling, explaining and describing at B1.
\end{grammarlens}

\subsection*{Your turn}
\begin{itemize}
\raggedright
  \item \emph{Say it:} \emph{otména}
  \item \emph{Say it:} \emph{otména}, once more
  \item \emph{Say it:} a sentence of your own with \emph{otména}
  \item \emph{Recall:} say the Russian for a delay, then the Russian for to delay, then say \emph{otména} again
\end{itemize}

\begin{tcolorbox}[breakable,colback=teal!4,colframe=teal!35!black,title={Before you move on}]
What does \ru{отмена} mean? (\textquotedblleft{}A cancellation\textquotedblright{}.) Say it once more.
\end{tcolorbox}
\section[\texorpdfstring{\ru{отменять} (otmenyát)}{отменять (otmenyát)}]{\ru{отменять} (otmenyát) — to cancel}

\label{lesson:RU-C423-otmenyat}



\begin{quote}
Before the new one: say the Russian for to delay, then the Russian for a cancellation.
\end{quote}

\subsection*{You'll want to know: \ru{отменять}}
\textbf{\ru{отменять}} (\emph{otmenyát}) — \textquotedblleft{}to cancel\textquotedblright{}.

\begin{quote}
\ru{В} \ru{плохую} \ru{погоду} \ru{часто} \ru{отменяют} \ru{рейсы}.
\end{quote}

\begin{quote}
\emph{V plokhúyu pogódu chásto otmenyáyut réysy.}
\end{quote}

\begin{quote}
In bad weather flights are often cancelled.
\end{quote}

Imperfective; perfective \emph{\ru{отменить}} (\emph{\ru{отменю}, \ru{отменишь}}).

\begin{cousinweb}
\emph{\ru{отменять}} is built from \emph{\ru{от}-} 'away' and \emph{\ru{менять}} 'to change', on \emph{\ru{мена}} 'an exchange', which is inherited from Proto-Slavic \emph{měna} (EtymDB).
\end{cousinweb}

\begin{grammarlens}[title={What you've built}]
One more word for telling, explaining and describing at B1.
\end{grammarlens}

\subsection*{Your turn}
\begin{itemize}
\raggedright
  \item \emph{Say it:} \emph{otmenyát}
  \item \emph{Say it:} \emph{otmenyát}, once more
  \item \emph{Say it:} a sentence of your own with \emph{otmenyát}
  \item \emph{Recall:} say the Russian for to delay, then the Russian for a cancellation, then say \emph{otmenyát} again
\end{itemize}

\begin{tcolorbox}[breakable,colback=teal!4,colframe=teal!35!black,title={Before you move on}]
What does \ru{отменять} mean? (\textquotedblleft{}To cancel\textquotedblright{}.) Say it once more.
\end{tcolorbox}
\section[\texorpdfstring{\ru{компенсация} (kompensátsiya)}{компенсация (kompensátsiya)}]{\ru{компенсация} (kompensátsiya) — compensation}

\label{lesson:RU-C423-kompensatsiya}



\begin{quote}
Before the new one: say the Russian for a cancellation, then the Russian for to cancel.
\end{quote}

\subsection*{You'll want to know: \ru{компенсация}}
\textbf{\ru{компенсация}} (\emph{kompensátsiya}) — \textquotedblleft{}compensation\textquotedblright{}.

\begin{quote}
\ru{За} \ru{задержку} \ru{пассажиры} \ru{получили} \ru{компенсацию}.
\end{quote}

\begin{quote}
\emph{Za zadérzhku passazhíry poluchíli kompensátsiyu.}
\end{quote}

\begin{quote}
The passengers got compensation for the delay.
\end{quote}

Feminine; formal.

\begin{cousinweb}
\emph{\ru{компенсация}} is usually traced through French \emph{compensation} to Latin \emph{compensātiō} 'a balancing' (EtymDB lists both), from \emph{compensāre} 'to weigh one thing against another' (Lewis \& Short).
\end{cousinweb}

\begin{grammarlens}[title={What you've built}]
One more word for telling, explaining and describing at B1.
\end{grammarlens}

\subsection*{Your turn}
\begin{itemize}
\raggedright
  \item \emph{Say it:} \emph{kompensátsiya}
  \item \emph{Say it:} \emph{kompensátsiya}, once more
  \item \emph{Say it:} a sentence of your own with \emph{kompensátsiya}
  \item \emph{Recall:} say the Russian for a cancellation, then the Russian for to cancel, then say \emph{kompensátsiya} again
\end{itemize}

\begin{tcolorbox}[breakable,colback=teal!4,colframe=teal!35!black,title={Before you move on}]
What does \ru{компенсация} mean? (\textquotedblleft{}Compensation\textquotedblright{}.) Say it once more.
\end{tcolorbox}
